% généré depuis source/12_bibliographie.md — ne pas modifier à la main

\chapter*{Bibliographie}
\addcontentsline{toc}{chapter}{Bibliographie}
\frontmark{Bibliographie}
\begin{multicols}{2}
\small\setstretch{1.15}\sloppy
\begingroup
\renewcommand{\chapter}[2]{}% le titre est déjà imprimé
\begin{thebibliography}{99}
\setlength{\itemsep}{2pt}
\bibitem{saxena2008data} A. Saxena and K. Goebel, ``Turbofan engine degradation simulation data set,'' NASA Ames Prognostics Data Repository, NASA Ames Research Center, Moffett Field, CA, 2008.
\bibitem{saxena2008damage} A. Saxena, K. Goebel, D. Simon, and N. Eklund, ``Damage propagation modeling for aircraft engine run-to-failure simulation,'' in \emph{Proc. Int. Conf. Prognostics and Health Management (PHM)}, 2008.
\bibitem{nf13306} AFNOR, \emph{Maintenance — Terminologie de la maintenance}, NF EN 13306, 2018.
\bibitem{jardine2006} A. K. S. Jardine, D. Lin, and D. Banjevic, ``A review on machinery diagnostics and prognostics implementing condition-based maintenance,'' \emph{Mech. Syst. Signal Process.}, vol. 20, no. 7, pp. 1483–1510, 2006.
\bibitem{benhmida} F. Ben Hmida, ``Estimation des indices FMDS,'' Support de cours, module « Sûreté de fonctionnement et maintenance », GMAM 3, ENSIT, Tunis, 2025.
\bibitem{lei2018} Y. Lei, N. Li, L. Guo, N. Li, T. Yan, and J. Lin, ``Machinery health prognostics: A systematic review from data acquisition to RUL prediction,'' \emph{Mech. Syst. Signal Process.}, vol. 104, pp. 799–834, 2018.
\bibitem{si2011} X.-S. Si, W. Wang, C.-H. Hu, and D.-H. Zhou, ``Remaining useful life estimation — A review on the statistical data driven approaches,'' \emph{Eur. J. Oper. Res.}, vol. 213, no. 1, pp. 1–14, 2011.
\bibitem{kaplanmeier1958} E. L. Kaplan and P. Meier, ``Nonparametric estimation from incomplete observations,'' \emph{J. Amer. Statist. Assoc.}, vol. 53, no. 282, pp. 457–481, 1958.
\bibitem{greenwood1926} M. Greenwood, ``The natural duration of cancer,'' His Majesty's Stationery Office, London, Reports on Public Health and Medical Subjects 33, 1926.
\bibitem{nelson1972} W. Nelson, ``Theory and applications of hazard plotting for censored failure data,'' \emph{Technometrics}, vol. 14, no. 4, pp. 945–966, 1972.
\bibitem{aalen1978} O. Aalen, ``Nonparametric inference for a family of counting processes,'' \emph{Ann. Statist.}, vol. 6, no. 4, pp. 701–726, 1978.
\bibitem{weibull1951} W. Weibull, ``A statistical distribution function of wide applicability,'' \emph{J. Appl. Mech.}, vol. 18, pp. 293–297, 1951.
\bibitem{bernard1953} A. Bernard and E. C. Bos-Levenbach, ``The plotting of observations on probability paper,'' \emph{Statistica Neerlandica}, vol. 7, pp. 163–173, 1953.
\bibitem{massey1951} F. J. Massey, ``The Kolmogorov-Smirnov test for goodness of fit,'' \emph{J. Amer. Statist. Assoc.}, vol. 46, no. 253, pp. 68–78, 1951.
\bibitem{akaike1974} H. Akaike, ``A new look at the statistical model identification,'' \emph{IEEE Trans. Autom. Control}, vol. 19, no. 6, pp. 716–723, 1974.
\bibitem{barlow1965} R. E. Barlow and F. Proschan, \emph{Mathematical Theory of Reliability}. New York, NY, USA: Wiley, 1965.
\bibitem{coble2009} J. Coble and J. W. Hines, ``Identifying optimal prognostic parameters from data: A genetic algorithms approach,'' in \emph{Proc. Annu. Conf. Prognostics and Health Management Society}, 2009.
\bibitem{wang2008} T. Wang, J. Yu, D. Siegel, and J. Lee, ``A similarity-based prognostics approach for remaining useful life estimation of engineered systems,'' in \emph{Proc. Int. Conf. Prognostics and Health Management (PHM)}, 2008.
\bibitem{gebraeel2005} N. Z. Gebraeel, M. A. Lawley, R. Li, and J. K. Ryan, ``Residual-life distributions from component degradation signals: A Bayesian approach,'' \emph{IIE Trans.}, vol. 37, no. 6, pp. 543–557, 2005.
\bibitem{si2012} X.-S. Si, W. Wang, C.-H. Hu, D.-H. Zhou, and M. G. Pecht, ``Remaining useful life estimation based on a nonlinear diffusion degradation process,'' \emph{IEEE Trans. Rel.}, vol. 61, no. 1, pp. 50–67, 2012.
\bibitem{heimes2008} F. O. Heimes, ``Recurrent neural networks for remaining useful life estimation,'' in \emph{Proc. Int. Conf. Prognostics and Health Management (PHM)}, 2008.
\bibitem{hochreiter1997} S. Hochreiter and J. Schmidhuber, ``Long short-term memory,'' \emph{Neural Comput.}, vol. 9, no. 8, pp. 1735–1780, 1997.
\bibitem{zheng2017} S. Zheng, K. Ristovski, A. Farahat, and C. Gupta, ``Long short-term memory network for remaining useful life estimation,'' in \emph{Proc. IEEE Int. Conf. Prognostics and Health Management (ICPHM)}, 2017.
\bibitem{babu2016} G. S. Babu, P. Zhao, and X.-L. Li, ``Deep convolutional neural network based regression approach for estimation of remaining useful life,'' in \emph{Database Systems for Advanced Applications (DASFAA)}, ser. Lecture Notes in Computer Science, vol. 9642. Springer, 2016.
\bibitem{li2018} X. Li, Q. Ding, and J.-Q. Sun, ``Remaining useful life estimation in prognostics using deep convolution neural networks,'' \emph{Reliab. Eng. Syst. Saf.}, vol. 172, pp. 1–11, 2018.
\bibitem{ramasso2014} E. Ramasso and A. Saxena, ``Performance benchmarking and analysis of prognostic methods for CMAPSS datasets,'' \emph{Int. J. Prognostics and Health Management}, vol. 5, no. 2, 2014.
\bibitem{davidsonpilon2019} C. Davidson-Pilon, ``lifelines: Survival analysis in Python,'' \emph{J. Open Source Softw.}, vol. 4, no. 40, pp. 1317, 2019.
\bibitem{pedregosa2011} F. Pedregosa et al., ``Scikit-learn: Machine learning in Python,'' \emph{J. Mach. Learn. Res.}, vol. 12, pp. 2825–2830, 2011.
\bibitem{friedman2001} J. H. Friedman, ``Greedy function approximation: A gradient boosting machine,'' \emph{Ann. Statist.}, vol. 29, no. 5, pp. 1189–1232, 2001.
\bibitem{ke2017} G. Ke, Q. Meng, T. Finley, T. Wang, W. Chen, W. Ma, Q. Ye, and T.-Y. Liu, ``LightGBM: A highly efficient gradient boosting decision tree,'' in \emph{Advances in Neural Information Processing Systems (NeurIPS)}, vol. 30, 2017.
\bibitem{cortes1995} C. Cortes and V. Vapnik, ``Support-vector networks,'' \emph{Mach. Learn.}, vol. 20, pp. 273–297, 1995.
\bibitem{platt1999} J. Platt, ``Probabilistic outputs for support vector machines and comparisons to regularized likelihood methods,'' in \emph{Advances in Large Margin Classifiers}. Cambridge, MA, USA: MIT Press, 1999.
\bibitem{kingma2015} D. P. Kingma and J. Ba, ``Adam: A method for stochastic optimization,'' in \emph{Proc. Int. Conf. Learning Representations (ICLR)}, 2015.
\bibitem{paszke2019} A. Paszke et al., ``PyTorch: An imperative style, high-performance deep learning library,'' in \emph{Advances in Neural Information Processing Systems (NeurIPS)}, vol. 32, 2019.
\bibitem{hanley1982} J. A. Hanley and B. J. McNeil, ``The meaning and use of the area under a receiver operating characteristic (ROC) curve,'' \emph{Radiology}, vol. 143, no. 1, pp. 29–36, 1982.
\bibitem{efron1979} B. Efron, ``Bootstrap methods: Another look at the jackknife,'' \emph{Ann. Statist.}, vol. 7, no. 1, pp. 1–26, 1979.
\bibitem{cox1972} D. R. Cox, ``Regression models and life-tables,'' \emph{J. Roy. Statist. Soc. B}, vol. 34, no. 2, pp. 187–220, 1972.
\bibitem{lundberg2017} S. M. Lundberg and S.-I. Lee, ``A unified approach to interpreting model predictions,'' in \emph{Advances in Neural Information Processing Systems (NeurIPS)}, vol. 30, 2017.
\bibitem{vaswani2017} A. Vaswani et al., ``Attention is all you need,'' in \emph{Advances in Neural Information Processing Systems (NeurIPS)}, vol. 30, 2017.
\bibitem{ishwaran2008} H. Ishwaran, U. B. Kogalur, E. H. Blackstone, and M. S. Lauer, ``Random survival forests,'' \emph{Ann. Appl. Statist.}, vol. 2, no. 3, pp. 841–860, 2008.
\bibitem{raissi2019} M. Raissi, P. Perdikaris, and G. E. Karniadakis, ``Physics-informed neural networks: A deep learning framework for solving forward and inverse problems involving nonlinear partial differential equations,'' \emph{J. Comput. Phys.}, vol. 378, pp. 686–707, 2019.
\bibitem{tao2019} F. Tao, H. Zhang, A. Liu, and A. Y. C. Nee, ``Digital twin in industry: State-of-the-art,'' \emph{IEEE Trans. Ind. Informat.}, vol. 15, no. 4, pp. 2405–2415, 2019.
\end{thebibliography}
\endgroup
\end{multicols}
